\section{Density of periodic points}


\subsection{General strategy}

  \paragraph{Historical note.} On \Yang{}

  First observe that if $X$ is defined over $\overline{\bbF_p}$, then for every point $x \in X(\overline{\bbF_p})$, $x$ is defined over some finite extension $\bbF_{p^n}$ of $\bbF_p$.
  In other words, the morphism $x: \Spec \overline{\bbF_p} \to X$ factors as $\Spec \overline{\bbF_p} \xrightarrow{\sigma} \Spec \bbF_{p^n} \xrightarrow{x_0} X$ for some $n > 0$.
  Applying the absolute Frobenius morphism $\Frob_{q}$, we get the diagram
  \[ \begin{tikzcd}
    x: & \Spec \overline{\bbF_p} \arrow[r, "\sigma"] & \Spec \bbF_{p^n} \arrow[r, "x_0"] & X,  \\
    \Frob_{q}(x): & \Spec \overline{\bbF_p} \arrow[r, "\sigma"] & \Spec \bbF_{p^n} \arrow[d,"\Frob_{q}"] \arrow[r, "x_0"] & X \arrow[d,"\Frob_{q}"] \\
      &  & \Spec \bbF_{p^n} \arrow[r, "x_0"] & X.
  \end{tikzcd} \]
  When $p^n = q$, the arrow $\Spec \bbF_{p^n} \xrightarrow{\Frob_q} \Spec \bbF_{p^n}$ becomes the identity, and hence $\Frob_q(x) = \Frob_p^n(x) = x$.
  Therefore, every point $x \in X(\overline{\bbF_p})$ is a periodic point of $\Frob_p$.
  
  For a general rational self-map $f: X \dashrightarrow X$ defined over $\overline{\bbF_p}$, we will try to find intersection points of the graph of $f$ and the graph of $\Frob_p^n$ for some $n > 0$, which will be periodic points of $f$.
  To find such points, we will use the Hrushovski-Lang-Weil estimate, which is a generalization of the classical Lang-Weil estimate for counting points on varieties over fields of positive characteristic.

  \begin{theorem}[{Hrushovski-Lang-Weil estimate, ref. \cite[Theorem 1.1]{Hrushovski2004elementary}}]\label{thm:Hrushovski-Lang-Weil_estimate}
    Let $\kkk$ be an algebraically closed field of characteristic $p > 0$ and $X$ be an affine variety over $\kkk$.
    Let $q$ be a power of $p$ and $\Frob_{q,/\kkk}: X \to X^{(q)}$ be the relative Frobenius morphism.
    Let $S \subset X \times X^{(q)}$ be a subvariety of the same dimension as $X$ such that 
    the two projections $p_1: S \to X$ and $p_2: S \to X^{(q)}$ are both dominant and $p_2$ is quasi-finite.
    Then there is a constant $c$ depending only on the degree of $S$ and $X$ such that for all $q > c$, we have 
    \[ \# (S \cap \Gamma_{\Frob_{q,/\kkk}})(\kkk) = \frac{\deg(p_1)}{\deg(p_2)_{\mathrm{insep}}} q^{\dim X} + O(q^{\dim X - \frac{1}{2}}), \]
    where $\Gamma_{\Frob_{q,/\kkk}}$ is the graph of $\Frob_{q,/\kkk}$ and $\deg(p_2)_{\mathrm{insep}}$ is the inseparable degree of $p_2$.
  \end{theorem}

  The statement of \cref{thm:Hrushovski-Lang-Weil_estimate} is stated in full generality but a little bit technical, and we will only use the following special case.

  \begin{lemma}[{ref. \cite[Corollary 1.2]{Hrushovski2004elementary}}]
    Let $X$ be an affine variety defined over $\bbF_q$ and $S \subset X \times X$ be a subvariety over $\overline{\bbF_q}$ such that the two projections $S \to X$ are both dominant.
    Then for any proper closed subset $W \subset X$, for large enough $m$, there exists a point $x \in X(\overline{\bbF_q})$ such that $(x, \Frob_q^m(x)) \in S$ and $x \notin W$.
  \end{lemma}

  \begin{lemma}
    Let $X$ be a variety over $\overline{\bbF_p}$ and $f: X \dashrightarrow X$ be a dominant rational self-map.
    Then the set of periodic points of $f$ is Zariski dense in $X$.
  \end{lemma}
  \begin{proof}
    Let $U$ be an non-empty affine open subset of $X$.
    We only need to show that there exists a periodic point of $f$ in $U$.
    Let $q$ be a power of $p$ such that $U$ and $f$ is defined over $\bbF_q$.
    Set $\Gamma_f$ to be the graph of $f$ in $U \times U$ and $\Gamma_{\Frob_q}$ to be the graph of $\Frob_q$ in $U \times U$.

    Choose $W \subset U$ to be a proper closed subset such that $f$ is defined on $U \setminus W$ and $f(U \setminus W) \subset U$.
    By the Hrushovski-Lang-Weil estimate, for large enough $m$, there exists a point $x \in U(\overline{\bbF_q})$ such that $(x, \Frob_q^m(x)) \in \Gamma_f$ and $x \notin W$.
    This means that $f(x) = \Frob_q^m(x)$.

    Since $X$ is defined over $\bbF_q$, there exists $n > 0$ such that $\Frob_q^{mn}(x) = x$.
    Note that $f$ is defined over $\bbF_q$, we have 
    \[ f^n(x) = f^{n-1}(f(x)) = f^{n-1}(\Frob_q^m(x)) = \Frob_q^m(f^{n-1}(x)) = \cdots = \Frob_q^{mn}(x) = x. \]
    We are done.
    % \Yang{}
  \end{proof}

  \begin{definition}
    We say a periodic point $x$ of $f$ is \emph{isolated} if there exists a period $r$ of $x$ such that $x$ is an isolated fixed point of $f^r$.
    % \Yang{}
  \end{definition}

  An important fact is that isolated periodic points can be lifted to generic fibers in a family, which allows us to reduce the density of periodic points to the density of isolated periodic points on special fibers.

  \begin{lemma}
    Let $S = \Spec R$ be spectrum of a discrete valuation ring with special point $s$ and generic point $\eta$. 
    Let $X$ be a scheme flat and of finite type over $S$, and $f: X \dashrightarrow X$ be a rational self-map over $S$.
    Assume that the special fiber $X_s$ is irreducible and the restriction $f_s: X_s \dashrightarrow X_s$ is a dominant rational self-map. 
    
    If $x\in X_s$ is an isolated periodic point of $f_s$ and $X_s$ is regular at $x$, then there exists a periodic point $y\in X_\eta$ of $f_\eta$ such that the closure of $y$ in $X$ contains $x$.
    % \Yang{}
  \end{lemma}
  \begin{proof}
    Replace $f$ by a suitable iterate, we may assume that $x$ is an isolated fixed point of $f_s$.
    Let $\Gamma_f$ be the graph of $f$ in $X \times_S X$ and $\Delta_X$ be the diagonal of $X \times_S X$.
    Since $X_s$ is regular at $x$, $X$ is flat over $S$ at $x$, and $S$ is regular, we have $X$ is regular at $x$.
    Then the embedding $\Delta_X \hookrightarrow X \times_S X$ is a regular embedding near $x$.
    By \Yang{}, we have 
    \[ \dim_x (\Gamma_f \cap \Delta_X) \geq \dim_x \Gamma_f + \dim_x \Delta_X - \dim_x (X \times_S X) = 1. \]
    Let $Z$ be an irreducible component of $\Gamma_f \cap \Delta_X$ containing $(x,x)$, then $\dim Z \geq 1$.
    Since $x$ is an isolated fixed point of $f_s$, the special fiber of $Z$ is just the point $x$.
    Then $Z \to S$ is dominant.
    In particular, the generic fiber of $Z$ over $S$ is non-empty.
    Let $y$ be a point in the generic fiber of $Z$, then $y$ is a periodic point of $f_\eta$ and the closure of $y$ in $X$ contains $x$.
    % \Yang{}
  \end{proof}

  \begin{example}
    When $x$ is a non-isolated periodic point of $f_s$, we may not be able to lift $x$ to a periodic point of $f_\eta$ in general.
    For example, let $S = \Spec \bbZ_p$ and $X = \bbA^1_S$.
    Consider the morphism $f: X \to X$ defined by $f(x) = x + p$.
    Over the special fiber, $f_s$ is the identity map, and every point is a periodic point.
    However, over the generic fiber, $f_\eta$ has no periodic point.

    We have $X \times_S X = \Spec \bbZ_p[x,y]$, $\Delta_X = \{x=y\} = \Spec \bbZ_p[x,y]/(x-y)$ and $\Gamma_f = \{y = x + p\} = \Spec \bbZ_p[x,y]/(y-x-p)$.
    Then 
    \[ \Delta_X \cap \Gamma_f = \Spec \bbZ_p[x,y]/(x-y, y-x-p) = \Spec \bbF_p[x].\]
    Its special fiber is $\Spec \bbF_p[x]$, which is one-dimensional, and its generic fiber is empty.
  \end{example}


\subsection{For cohomologically hyperbolic maps}

  \begin{proposition}
    Let $X$ be a variety and $f: X \dashrightarrow X$ be a rational self-map.
    Suppose that $f$ is cohomologically hyperbolic, then after replacing $X$ by an open dense subset, the set of periodic points of $f$ is isolated in $X$.
    \Yang{}
  \end{proposition}
  \begin{proof}
    Set $\mu_k = \mu_k(f)$ for $1 \leq k \leq d$ and $\mu_{d+1} = 0$ by convention.
    There exists $i$ such that $\mu_i > 1 > \mu_{i+1}$ by the definition of cohomological hyperbolicity.
    Choose $a,b > 0$ such that $\mu_i > a > 1 > b > \mu_{i+1}$ and $a \cdot b = \mu_i \cdot \mu_{i+1}$.
    Then there exists $\varepsilon > 0$ such that $\varepsilon^m \mu_i^m > a^m + b^m$ for all $m > m_1$ for some $m_1 > 0$.
    By \Yang{}, the divisor
    \[ (f^{2m})^*L - \varepsilon^m \mu_i^m (f^m)^*L + \mu_i^m \mu_{i+1}^m L \]
    is big for some ample line bundle $L$ and all $m > m_2$ for some $m_2 > 0$.
    Then fix a large enough $m > \max\{m_1, m_2\}$, and let 
    \[ M := (f^{2m})^*L - (a^m + b^m) (f^m)^*L + a^m b^m L. \]
    Then $M$ is big.
    In particular, $M$ is numerically equivalent to an effective $\bbR$-divisor $\sum a_j D_j$ with $a_j > 0$ and $D_j$ prime divisors.
    Set $Z = \bigcup D_j$ to be the support of $M$.

    Let $U = X \setminus Z$ be the complement of $Z$ in $X$, which is an open dense subset of $X$.
    We claim that every periodic point of $f$ in $U$ is isolated.
    Otherwise, there exists a curve $C$ intersecting $U$ such that $f^r|_C = \id_C$ for some $r > 0$.
    After iterating $f$, we may assume that $f(C) = C$.
    Since $f^n(C) = C$ intersects $U$ for all $n > 0$, we have \Yang{}
    \[ (f^n)^*M \cdot C = M \cdot f^n_*C \geq 0. \]
    Then we have 
    \[ (f^{(n+2)m})^*L \cdot C - (a^m + b^m) (f^{(n+1)m})^*L \cdot C + a^m b^m (f^{nm})^*L \cdot C \geq 0. \]
    Let $x_n = (f^{nm})^*L \cdot C$, then $x_n$ satisfies the linear recurrence relation
    \[ x_{n+2} - (a^m + b^m) x_{n+1} + a^m b^m x_n \geq 0. \]
    Since $b<1$, we have $x_1 - b x_0 = (f^m)^*L \cdot C - b L \cdot C = L \cdot f^m_*C - b L \cdot C = (1-b) L \cdot C > 0$.
    Hence $(f^{nm})^*L \cdot C = x_n >  a^{nm}$ for $n \gg 0$.\Yang{} 
    However, we have $(f^{nm})^*L \cdot C = L \cdot f^{nm}_*C = L \cdot C$ since $f^{nm}_*C = C$, which is a contradiction.

    \Yang{}
  \end{proof}

  \begin{theorem}
    Let $X$ be a variety and $f: X \dashrightarrow X$ be a rational self-map.
    If $f$ is cohomologically hyperbolic, then the set of periodic points of $f$ is Zariski dense in $X$.
  \end{theorem}
  \begin{proof}
    \Yang{}
  \end{proof}